% year: תשע"ז
% type: מבחן
% semester: ב
% moed: א
% source: exams/מבחנים/תשעז/מועד א תשעז סמסטר ב.pdf
% number: 2
% points: 15
% categories: ivt, func-analysis

\begin{question}
האם לפונקציה $f(x)=5-x-x^2-|x+3|$ יש מקסימום בקטע $[-10,10]$? אם כן, מצאו אותו.
\end{question}

\begin{hint}
$f$ רציפה בקטע סגור -- איזה משפט מבטיח קיום מקסימום? כדי למצוא אותו, פתחו את הערך המוחלט ובדקו כל אחד משני התחומים בנפרד.
\end{hint}

\begin{solution}
\textbf{קיום:} $f$ רציפה ב-$[-10,10]$ (סכום של פולינום ושל $-|x+3|$, שהיא רציפה). לפי משפט ויירשטראס, פונקציה רציפה בקטע סגור וחסום מקבלת בו מקסימום (ומינימום). לכן \textbf{כן}, יש מקסימום.

\textbf{מציאתו:} נפתח את הערך המוחלט.
\begin{itemize}
  \item עבור $-3\le x\le10$: $|x+3|=x+3$ ולכן $f(x)=5-x-x^2-x-3=2-2x-x^2=3-(x+1)^2$. זו פרבולה הפוכה עם קודקוד ב-$x=-1\in[-3,10]$, ולכן בתחום זה $f(x)\le3$ ושוויון רק ב-$x=-1$.
  \item עבור $-10\le x<-3$: $|x+3|=-(x+3)$ ולכן $f(x)=5-x-x^2+x+3=8-x^2$. כאן $x^2>9$, ולכן $f(x)<8-9=-1<3$.
\end{itemize}
(לחלופין: $f$ גזירה פרט ל-$x=-3$; $f'(x)=-2-2x$ עבור $x>-3$ מתאפסת ב-$x=-1$, ו-$f'(x)=-2x>0$ עבור $x<-3$; משווים את ערכי $f$ בנקודות החשודות $-10,-3,-1,10$: $f(-10)=-92$, $f(-3)=-1$, $f(-1)=3$, $f(10)=-118$.)

לכן המקסימום של $f$ בקטע הוא $\boxed{f(-1)=3}$, והוא מתקבל בנקודה $x=-1$.
\end{solution}
